
Converting transformers to sub-quadratic architectures such as Mamba achieves O(n) inference complexity but typically requires more gradient steps to adapt to new tasks compared to the original transformer. Standard conversion approaches treat architecture change and downstream adaptation as sequential stages: first convert the model by optimizing for output matching, then fine-tune for specific tasks. This work examines whether integrating task conditioning into the conversion process itself can preserve adaptation capability.

The motivation stems from an observation that knowledge distillation preserves outputs but not necessarily the internal structure that enables rapid adaptation. When a transformer is compressed into a sub-quadratic form, task-relevant functional representations may degrade. Post-hoc fine-tuning operates on these potentially degraded representations.

We hypothesize that task-relevant structure already exists in pretrained transformer hidden states and can be discovered without supervision. If such structure can be extracted into task embeddings and used to modulate state space dynamics during conversion training, the resulting model may retain adaptation capability.

This paper introduces Task-Conditioned Selective State Space (TC-SSM), which makes three contributions:

\begin{enumerate}
\item We demonstrate that K-means clustering discovers task structure in BERT hidden states with purity 0.74 (versus 0.20 random baseline for 5 tasks), enabling task conditioning without labeled data.
\end{enumerate}

\begin{enumerate}
\item We show that rank-32 low-rank projections modulating Mamba's $\Delta$, B, and C matrices based on task embeddings add only 1.01x parameters and achieve 0.85x inference time relative to vanilla Mamba.
\end{enumerate}

\begin{enumerate}
\item We validate on SuperGLUE classification tasks that TC-SSM achieves 0.25\% accuracy gap versus transformer baseline while adapting in 14 gradient steps.
\end{enumerate}

